\documentclass{article}
\usepackage{amsmath} 
\usepackage{graphicx}

\title{Copy}
\author{Vikas}
\date{\today}
\newcommand{\R}{\mathcal{R}}
\newcommand{\sq}[1]{#1^2}
\newtheorem{theorem}{Theorem}[section]
\newtheorem{lemma}{Lemma}

\begin{document}
\maketitle

\section{Introduction}
Hello, this is my first latex practical.
\\ \\
This is the first paragraph.
\\ \\
This is the second paragraph.
\\ \\
This is line one.
\\ \\
This is line two.
\\ \\
"This is quoted text."
\\ \\
 Dashes : 
 - % hyphen
 -- % en dash
 --- % em dash
\\ \\
Hello Mr.\ Buttler % No Space after period 
\\ \\
Special Symbols: \#  \$  \%  \&  \_  \{  \}  \~
\\ \\
Example : 
I scored 96.2\% in my class X boards.
\\ \\
 \today        % current date
 \\ \\
 \LaTeX        % LaTeX logo 
 \\ \\
 \TeX          % TeX logo
\\ \\
Text Styles : 
\textit{Hello}       % italic
(Italic)
\textbf{Hello}      % bold
(Bold)
\underline{Hello}   % underline
(Underline)
\\ \\
\textit{hermeneutics}  % italic
\footnote{Hermeneutics: the theory and methodology of interpretation, especially of biblical texts, wisdom literature, and philosophical texts.}
\section{Introduction}
This is the Introduction Section.
\begin{quote}
    This is a short quotation.
\end{quote}

\subsection{Sub Section : List}
Unordered List 
\\ \\ 
Stationary
\begin{itemize}
    \item Pencil
    \item Pen
    \item Highlighter
    \item Scale
\end{itemize}
Ordered List
\begin{enumerate}
    \item Mathew
    \item Stephen
    \item George
\end{enumerate}
Description List
\begin{description}
    \item[Pen] A Stationary
    \item[Stephen] A Name
\end{description}

\subsection{Displaying Formulas}
\[
[a + b]^2 =a^2 + b^2 + 2ab 
\]
\[
u^2 + v^2 = 2at 
\]
\section{Introduction}
Running LaTeX, Changing the type style, Accents, Symbols, Subscripts and Superscripts, Fractions, Roots, Ellipsis. 
\subsection{Type Style}
\textbf{Latex},
\textit{Latex},
\texttt{Latex},
\textsc{Latex}

\subsection{Accents}
\'a , \^a, \"a, \~a, \`a, \=a, \.a,\u{a}, \v{a}

\subsection{Symbols}
\#  \$  \%  \&  \_  \{  \}
\\ \\
$\alpha, \beta, \gamma, \leq, \geq, \neq, \infty, \partial, \nabla,\to$

\subsection{Subscripts and Superscripts}
Subscript: $a_1, b_2$,
Superscript: $x^2,y^2$,
Both: $n_{ij}^2$

\subsection{Fractions}
$\frac{a}{b}$,
$\frac{x+1}{y-1}$

\subsection{Roots}
$\sqrt{x}$,
$\sqrt[n]{x}$,
$\sqrt{a^2+b^2}$

\subsection{Ellipsis}
$1,2,3,\ldots,n$
\\ \\
$a_1+b_2+\cdots+a_n$
\section{Introduction}
Mathematical Symbols, Greek letters, Calligraphic letters, Log-like functions, Arrays, The array environment, Vertical alignment, Delimiters, Multiline formulas. 

\subsection{Mathematical Symbols}
$\leq, \geq,\neq, \approx, \equiv, \in, \notin, \subset, \subseteq, \cup, \cap, \forall, \exists, \rightarrow, \Rightarrow, \infty, \partial, \nabla$

\subsection{Greek Letters}
$\alpha, \beta, \gamma, \delta, \epsilon,\lambda, \mu, \pi, \sigma, \theta, \omega, \Gamma, \Delta$

\subsection{Calligraphic Letters}
$\mathcal{L} = \mathcal{F}(x)$

\subsection{Log-like Functions}
$\log x,\ln x,\sin x,\cos x,\tan x,\exp x,\max, \min$

\subsection{Arrays}
\[
\begin{array}{ccc}
    1 & 2 & 3 \\
    3 & 4  & 4\\
    5 & 6 & 5
\end{array}
\]
\subsection{The Array Environment}
\[
\left(
\begin{array}{cc}
    1 & 0 \\
    0 & 1
\end{array}
\right)
\]
\subsection{Vertical Alignment}
\[
\begin{array}{ccc}
   x+y  & =& 1 \\
   x-y  & =& 2
\end{array}
\]

\subsection{Delimiters}
\[
\left(\frac{a}{b} \right)
\]
\[
\left( \frac{x+1}{y-1} \right) 
\]

\subsection{Multiline Formulas}
\begin{align}
    x+y&=1\\
    x-y&=2
\end{align}

\begin{align}
    x+y&=1\\
    x-y&=2
\end{align}
\section{Introduction}
Putting one thing above another, Over and underlining, Accents, Stacking symbols, Spacing in math mode, Changing style in math mode, Type style, Math style. 

\subsection{Putting One Thing Above Another}
\[
\overset{def}{=} \quad \underset{n \to \infty}{\lim}
\]

\subsection{Over and Underlining}
\[
\overline{z}=x-iy
\]
\[
\underline{Research Methodology}
\]

\subsection{Accents}
\[
\hat{x} + \vec{y}
\]

\subsection{Stacking Symbols}
\[
A \stackrel{def}{=} B
\]

\subsection{Spacing in math mode}
\[
a\,b \quad a\;b \quad a\!b
\]

\subsection{Changing Style in Math Mode}
\[
{\displaystyle \frac{a}{b}} \quad {\textstyle \frac{a}{b}}
\]

\subsection{Type Style}
\[
\mathbf{v}, \mathcal{F}
\]

\subsection{Math Style}
\[
\sum_{i=1}^{n} x_i \quad 
{\scriptstyle \sum_{i=1}^{n} x_i}
\]
\section{Introduction}
Defining commands, Defining environments, Theorems.

\subsection{Defining Commands}
Set of Real Numbers: $\mathcal{R}$
\\ \\
$\sqrt{x} + \sq{y}$

\subsection{Theorems}
\begin{theorem}
    Every non-empty subset of 
 that is bounded above has a supremum (least upper bound) in $\R$
\end{theorem}
\begin{theorem}
    Between any two real numbers, there exists a rational number.
\end{theorem}
\begin{theorem}
    All even numbers are divisible by 2.
\end{theorem}
\begin{lemma}
    2 divides 10.
\end{lemma}

\subsection{Algebra}
\begin{theorem}
    $a+b=b+a$
\end{theorem}


\section{Introduction}

\subsection{Figures and Tables}

\begin{figure}[h]
    \centering
    \includegraphics[width=0.3\linewidth]{vinay.jpg}
    \caption{Sample Figure}
    \label{fig:sample}
\end{figure}

\begin{table}[h]
    \centering
    \begin{tabular}{c|c}
        A & B \\
        1 & 2
    \end{tabular}
    \caption{Sample Table}
    \label{tab:sample}
\end{table}

\subsection{Marginal Notes}
This is a sample text.\marginpar{Note in margin}

\subsection{Tabbing Environment}
\begin{tabbing}
    Name~~~ \= Age~~~ \= City \\
    Ram \> 20 \> Delhi \\
    Shyam \> 22 \> Mumbai \\
\end{tabbing}

\subsection{Tabular Environment}
\begin{tabular}{c|c|c}
    Name & Age & City \\
    Ram & 20 & Delhi \\
    Shyam & 22 & Mumbai
\end{tabular}

\subsection{Cross-References}
\label{sec:intro}
See Section~\ref{tab:sample}
\newpage

\section{Bibliography and Citation}

A ring of endomorphisms of a module plays a very important role in many parts of mathematics \cite{nobusawa1964generalization}. As a generalization of this idea, we can consider homomorphisms between modules \cite{nobusawa1964generalization}. 

Further developments in ring theory and related problems were studied extensively in later works \cite{kaplansky1970problems}.

We can also cite multiple references together when needed \cite{nobusawa1964generalization, kaplansky1970problems}.

\bibliographystyle{unsrt}
\bibliography{bibfiles}

\end{document}
